% title:          Functional equation in x plus 2y
% area:           limits-continuity, functional-equations
% methods:        substitution
% difficulty:     easy
% difficulty-ai:  easy
% status:
% review:         none
% --- history: all optional, leave EMPTY if unknown ---
% origin:         jarnik
% origin-date:    2017-03-31
% origin-number:  Category I, Problem 1
% also-in:
% references:     Earliest known appearance (to the best of our knowledge): Vojtech Jarnik International Mathematical Competition, 27th edition, Ostrava, 31 March 2017, Category I, Problem 1 - https://vjimc.osu.cz/storage/uploads/j27problems1.pdf
% history:        ai
\begin{problem}
    Let \( f \in C^0\left(\mathbb{R},\mathbb{R}\right) \) be such that, for all \( x, y \in
    \mathbb{R} \),
    \[
        f\left(x+2y\right) = 2f\left(x\right)f\left(y\right).
    \]
    Show that \( f \) is constant.
\end{problem}
\begin{solution}
    \textit{Solution 1.}
    \\\\
    For all \( x, y \in \mathbb{R} \), we have
    \[
        f\left(x+2y\right) = 2f\left(x\right)f\left(y\right) = 2f\left(y\right)f\left(x\right) =
        f\left(y+2x\right).
    \]
    Now, as the matrix \( \begin{pmatrix} 1 & 2 \\ 2 & 1 \end{pmatrix} \in
    \mathbb{R}^{2 \times 2} \) has full rank (its determinant is \( 1 \cdot 1-2 \cdot 2 = -3 \neq
    0 \)), for all \( a, b \in \mathbb{R} \), there exist \( x, y \in \mathbb{R} \) with
    \[
        \begin{pmatrix} 1 & 2 \\ 2 & 1 \end{pmatrix}\begin{pmatrix} x \\ y \end{pmatrix} =
        \begin{pmatrix} a \\ b \end{pmatrix},
    \]
    i.e., \( x+2y = a \) and \( y+2x = b \), so that \( f\left(a\right) = f\left(x+2y\right) =
    f\left(y+2x\right) = f\left(b\right) \). Hence, \( f \) must be constant.
    \begin{remark}
        This solution does not use the continuity of \( f \).
    \end{remark}
    \begin{remark}
        Due to the symmetry of \( 2f\left(x\right)f\left(y\right) \) in \( x \) and \( y \),
        Solution 1 only uses that \( f\left(u\left(v\right)\right) =
        f\left(w\left(v\right)\right) \) for all \( v \in \mathbb{R}^2 \), where \( u, w \colon
        \mathbb{R}^2 \to \mathbb{R} \) are the linear forms \( \left(x, y\right) \mapsto x+2y \) and
        \( \left(x, y\right) \mapsto y+2x \), which are linearly independent. Hence the following
        generalisation, with the same proof: let \( K \) be a field, \( m \in
        \mathbb{N}_{\geq 2} \), \( u, w \colon K^m \to K \) two linearly independent \( K \)-linear
        forms, \( Z \) a set and \( f \colon K \to Z \) a map such that
        \( f\left(u\left(v\right)\right) = f\left(w\left(v\right)\right) \) for all \( v \in K^m \);
        then \( f \) is constant. Indeed, the \( K \)-linear map \( K^m \to K^2 \), \( v \mapsto
        \left(u\left(v\right), w\left(v\right)\right) \) has rank \( 2 \) (its matrix has the two
        rows \( u \) and \( w \), which are linearly independent, and the rank of a matrix is
        standardly equal to the dimension of the span of its rows), hence is surjective, so that,
        for all \( a, b \in K \), there exists \( v \in K^m \) with \( u\left(v\right) = a \) and
        \( w\left(v\right) = b \), and then \( f\left(a\right) = f\left(u\left(v\right)\right) =
        f\left(w\left(v\right)\right) = f\left(b\right) \). In its most general form, the argument
        shows: if \( X \), \( Y \) and \( Z \) are sets and \( g, h \colon X \to Y \) and
        \( f \colon Y \to Z \) are maps such that \( f \circ g = f \circ h \) and the map
        \( X \to Y \times Y \), \( v \mapsto \left(g\left(v\right), h\left(v\right)\right) \) is
        surjective, then \( f \) is constant.
    \end{remark}
    \noindent\textit{Solution 2.}
    \\\\
    For all \( x \in \mathbb{R} \), we have \( f\left(x\right) = f\left(x+2 \cdot 0\right) =
    2f\left(x\right)f\left(0\right) \), so that \( f\left(x\right)\left(1-2f\left(0\right)\right) =
    0 \). If \( 1-2f\left(0\right) \neq 0 \), then \( f\left(x\right) = 0 \) for all \( x \in
    \mathbb{R} \), and hence \( f \) is the constant function \( 0 \).
    \\\\
    Otherwise, \( 1-2f\left(0\right) = 0 \), so \( f\left(0\right) = \frac{1}{2} \), and we must
    have, for all \( x \in \mathbb{R} \),
    \[
        f\left(2x\right) = f\left(0+2x\right) = 2f\left(0\right)f\left(x\right) = f\left(x\right).
    \]
    In particular, we easily have by induction, for all \( n \in \mathbb{Z} \) and all \( x \in
    \mathbb{R} \), \( f\left(x\right) = f\left(x \cdot 2^n\right) \). Fix \( x \in \mathbb{R} \).
    Since \( x \cdot 2^{-n} \overset{n \to +\infty}{\longrightarrow} 0 \) (as \( 2^{-n} =
    \frac{1}{2^n} \overset{n \to +\infty}{\longrightarrow} 0 \)) and \( f \) is continuous, we get
    \[
        f\left(x\right) = f\left(x \cdot 2^{-n}\right) \overset{n \to +\infty}{\longrightarrow}
        f\left(0\right),
    \]
    where the equality was seen above (with \( -n \) in place of \( n \)). As the left-hand side
    does not depend on \( n \), we have \( f\left(x\right) = f\left(0\right) \), and hence \( f \)
    is the constant function \( f\left(0\right) \).
    \\\\
    So, in both cases, \( f \) is constant.

    % Solution 3 is a student's solution from the seminar repository (compmath,
    % solutions/sol_continuity.tex, lines 62-70), copied word for word on 2026-10-22 at Antoine's
    % request (only the math delimiters, and the source file's own macros and environments, changed
    % so that it compiles here); not reviewed yet.
    \textit{Solution 3.}
    \( x=-z, y=z \) yields \( f(z)=2f(-z)f(z) \) \\
    \( x=z, y=-z \) yields \( f(-z) = 2f(-z)f(z) \) \\
    Combining these yields \( f(-z)=2f(z)f(-z)=f(z) \) \\
    In particular, using the \( x=-z, y=z \) substitution again, we get \( f(z)=2f(z)^2 \forall z
    \in \mathbb{R} \). \\
    We will now exclude the case where \( f(x)=0 \) for some \( x \) and \( f(x)=\frac{1}{2} \) for
    others, without using continuity (as it is not needed). \\
    Indeed suppose there exists some \( x_0 \in \mathbb{R} \) such that \( f(x_0)=0 \). In this
    case, use \( y=x_0, x=z-2x_0 \) to get \( f(z)=0 \forall z \in \mathbb{R} \). Hence we either do
    not have such an \( x \), in which case \( f(x)=\frac{1}{2} \forall x \in \mathbb{R} \) or
    \( f(x)=0 \forall x \in \mathbb{R} \)

\end{solution}
